\documentclass[12pt]{article}
\usepackage[utf8]{inputenc}
\usepackage[T1]{fontenc}
\usepackage{hyperref}
\title{Kokoda Language Resources for K-12 Curriculum Development}
\author{Compiled by Claude Sonnet 5.5 (claude-sonnet-5-5) for Andr\'as Kornai}
\date{}
\begin{document}
\maketitle

\section{Introduction}
Kokoda (ISO 639-3 xod; Emeyode) is a Papuan language of the Bird's Head, spoken in Kokoda District, South Sorong Regency, in Indonesia's Southwest Papua. Sources disagree on the number of dialects (three or four, including Yamueti). A 2026 study of the Warmon community reports a significant decline in active regional-language skills among Kokoda youth. No school curriculum or primer was found in the searches.

\section{Primary Educational Resources}
None found.

\section{Linguistic and Reference Materials}
\begin{itemize}
\item Lourens de Vries, ``The Kokoda language'', in \emph{A Short Grammar of Inanwatan} (Pacific Linguistics 560, 2004, pp. 130--137), a short grammar sketch. The text was not retrieved. Bibliographic record: \url{https://glottolog.org/resource/reference/id/29954}
\item Uppsala University, ``Two talks on South Bird's Head languages'' (2026): a talk by Konstantin Henke on the Kokoda verb system, based on preliminary documentation: \url{https://www.uu.se/institution/lingvistik-och-filologi/kalendarium/arkiv/2026-06-05-two-talks-on-south-birds-head-languages}
\item Wikipedia, ``Kokoda language'': \url{https://en.wikipedia.org/wiki/Kokoda_language}
\item Wikipedia, ``South Sorong Regency'' (regional context): \url{https://en.wikipedia.org/wiki/South_Sorong_Regency}
\end{itemize}

\section{Religious and Community Resources}
\begin{itemize}
\item Result titled ``301 xod 001'' on bakat.sabda.org, which appears to be a Christian audio or resource catalogue entry for xod. Its content was not examined, so its nature is unconfirmed: \url{https://bakat.sabda.org/301-xod-001.html}
\end{itemize}

\section{Digital and Audio Resources}
See the bakat.sabda.org entry above, which may be relevant (unconfirmed).

\section{Curriculum Development Context}
A sociolinguistic study of the Kokoda community and the use of its language in preserving identity is available (Indonesian): \url{https://e-jurnal.iainsorong.ac.id/index.php/El-Tarbawi/article/view/3136}. It documents language shift among youth in the Sorong area. With documentation still preliminary, a curriculum would depend on collaboration with Kokoda speakers and the documentation team. The Indonesian local-content (muatan lokal) framework could be a vehicle, but no existing school materials were found.
\end{document}